% answers_short.tex —— 2007 年·普通高等学校招生考试·江苏卷 · 简短答案速查
% 由 tools/make_short_answers.py 自 answers.tex 的 \gksolhead/\gkanswer 提取，
% 并与 2007/2007_数学真题_江苏_答案要点.md 逐题比对后生成，勿手改（改源 MD 后重跑）。

\begin{gknotice}
\textbf{使用说明}：本节只给最终答案便于快速核对，完整推导见解析版。答案为编写者独立推导，\textbf{非官方标准答案}。
\end{gknotice}

\gkpart{选择题}
\begin{center}\begin{tabular}{ccccccccccc}
\toprule
题号 & 1 & 2 & 3 & 4 & 5 & 6 & 7 & 8 & 9 & 10 \\
\midrule
答案 & D & A & A & C & D & B & B & A & C & B \\
\bottomrule
\end{tabular}\end{center}

\gkpart{填空题}
\begin{itemize}[leftmargin=2.6em]
  \item[\textbf{11.}] $\dfrac12$。
  \item[\textbf{12.}] $75$。
  \item[\textbf{13.}] $32$。
  \item[\textbf{14.}] $\dfrac{6\sqrt5}{5}$。
  \item[\textbf{15.}] $\dfrac54$。
  \item[\textbf{16.}] $10\sin\dfrac{\pi t}{60}$（$0\leqslant t\leqslant60$）。
\end{itemize}

\gkpart{解答题}
\begin{itemize}[leftmargin=2.6em]
  \item[\textbf{17.}] （1）$0.05$；（2）$0.99$；（3）$0.02$。
  \item[\textbf{18.}] （1）（2）证明见解析；（3）$\tan\theta=\sqrt{13}$。
  \item[\textbf{19.}] （1）$c=2$；（2）证明见解析；（3）逆命题成立。
  \item[\textbf{20.}] （1）（2）证明见解析；（3）存在，$q=\dfrac{\sqrt5-1}{2}$。
  \item[\textbf{21.}] （1）$d=0$；（2）$0\leqslant c<4$；（3）$0\leqslant c\leqslant4$。
\end{itemize}
